% titel: Schwebung
% skript: Kap. Schwingungen

% Schwebung zweier nahezu gleicher Frequenzen; Schwebungsfrequenz = Differenz.

Zwei Töne mit nahe beieinanderliegenden Frequenzen erzeugen ein periodisches
Lauter-/Leiserwerden --- die Schwebung. Die Zeit zwischen zwei Lautstärke-Maxima
ist die halbe Schwebungsperiode.

\begin{werkzeugkasten}[Werkzeuge]
\textbf{Überlagerung} zweier Kosinus-Schwingungen gleicher Amplitude, nahe
Frequenzen $f_1,f_2$ (mit Mittel $f_m=(f_1+f_2)/2$ und Differenz
$\Delta f=|f_1-f_2|$):
\begin{equation}
  \cos(2\pi f_1t)+\cos(2\pi f_2t)
  =2\cos(2\pi f_m t)\cos(\pi\,\Delta f\, t)\;.
\end{equation}

\textbf{Schwebungsfrequenz} $=$ Frequenz, mit der die \emph{Lautstärke}
(Envelope $|\cos(\pi\Delta f\,t)|$) maximal wird:
\begin{equation}
  f_{\mathrm{Schwebung}}=\Delta f=|f_1-f_2|\;,
  \qquad
  T_{\mathrm{Maximum}}=\frac{1}{\Delta f}\;.
\end{equation}
Der Envelope hat die Periode $2/\Delta f$, die beiden Maxima pro Envelope
liegen also $\Delta f^{-1}$ auseinander.

\textbf{Einheiten:} Frequenzen in \unit{\hertz}, Zeiten in \unit{\second}.
\end{werkzeugkasten}

\begin{werkzeugkasten}[Vorgehen]
\begin{itemize}
  \item Aus dem Abstand $T_{\max}$ zweier Lautstärke-Maxima direkt
       $\Delta f=1/T_{\max}$.
  \item Nicht mit der Mittelfrequenz verwechseln: Die Töne selbst schwingen
       mit $\approx f_m$, hörbar wird nur die Envelope-Frequenz $\Delta f$.
\end{itemize}
\end{werkzeugkasten}

\begin{achtung}[Typische Fallen]
\begin{itemize}
  \item $T_{\max}$ ist der Abstand zweier Maxima, also $\Delta f^{-1}$
       (nicht $2\Delta f^{-1}$, das wäre die Envelope-Periode).
  \item $\Delta f$ ist stets die \emph{Differenz} der beiden Frequenzen,
       unabhängig von der Tonhöhe.
\end{itemize}
\end{achtung}
